% =========================================================================
% DOCUMENTCLASS
% =========================================================================
\documentclass[12pt,oneside]{book}

% =========================================================================
% METADATOS
% =========================================================================
\newcommand{\universidad}{Universidad de Buenos Aires}
\newcommand{\carrera}{Carrera de Arquitectura}
\newcommand{\materia}{Legislación y Práctica Profesional}
\newcommand{\titulo}{El proceso judicial: estrategia, demanda y prueba}
\newcommand{\autor}{Isaac Edgar Camacho Ocampo}
\newcommand{\anio}{2026}

% =========================================================================
% PAQUETES
% =========================================================================
\usepackage[utf8]{inputenc}
\usepackage[T1]{fontenc}
\usepackage[spanish,es-nodecimaldot]{babel}

\usepackage[
    top=2.5cm,
    bottom=2.5cm,
    left=3cm,
    right=2.5cm,
    headheight=15pt
]{geometry}

\usepackage{amsmath}
\usepackage{amssymb}
\usepackage{mathtools}

\usepackage{graphicx}
\usepackage{float}
\usepackage{caption}
\usepackage{subcaption}

\usepackage{booktabs}
\usepackage{array}
\usepackage{longtable}
\usepackage{tabularx}

\usepackage{enumitem}

\usepackage{xcolor}
\usepackage[most]{tcolorbox}

\usepackage{fancyhdr}
\usepackage{ragged2e}

\usepackage[
    colorlinks=true,
    linkcolor=blue!50!black,
    citecolor=green!40!black,
    urlcolor=blue!60!black,
    pdftitle={\titulo},
    pdfauthor={\autor},
    pdfsubject={Apuntes universitarios},
    bookmarks=true
]{hyperref}

% =========================================================================
% CONFIGURACIÓN
% =========================================================================
\graphicspath{
    {./img/}
    {./graficos/}
    {./figuras/}
}

\setcounter{tocdepth}{2}

\setlength{\parindent}{0pt}
\setlength{\parskip}{0.5em}

\setlist[itemize]{
    topsep=0.3em,
    itemsep=0.2em
}

\setlist[enumerate]{
    topsep=0.3em,
    itemsep=0.2em
}

\clubpenalty=10000
\widowpenalty=10000

% =========================================================================
% COMANDOS
% =========================================================================
\newcommand{\abs}[1]{\left|#1\right|}
\newcommand{\norm}[1]{\left\lVert#1\right\rVert}
\newcommand{\vect}[1]{\mathbf{#1}}
\newcommand{\deriv}[2]{\frac{d#1}{d#2}}
\newcommand{\pderiv}[2]{\frac{\partial#1}{\partial#2}}

% =========================================================================
% ENTORNOS / CAJAS
% =========================================================================
\newtcolorbox{definition}[1]{
    enhanced,
    breakable,
    title={Definición: #1},
    fonttitle=\bfseries,
    colback=blue!3,
    colframe=blue!50!black,
    boxrule=0.8pt,
    arc=2mm
}

\newtcolorbox{important}{
    enhanced,
    breakable,
    title={Importante},
    fonttitle=\bfseries,
    colback=yellow!8,
    colframe=orange!70!black,
    boxrule=0.8pt,
    arc=2mm
}

\newtcolorbox{example}[1]{
    enhanced,
    breakable,
    title={Ejemplo: #1},
    fonttitle=\bfseries,
    colback=green!4,
    colframe=green!40!black,
    boxrule=0.8pt,
    arc=2mm
}

\newtcolorbox{formula}[1]{
    enhanced,
    breakable,
    title={Fórmula / Regla: #1},
    fonttitle=\bfseries,
    colback=purple!4,
    colframe=purple!50!black,
    boxrule=0.8pt,
    arc=2mm
}

\newtcolorbox{exercise}[1]{
    enhanced,
    breakable,
    title={Ejercicio: #1},
    fonttitle=\bfseries,
    colback=gray!5,
    colframe=gray!60!black,
    boxrule=0.8pt,
    arc=2mm
}

\newtcolorbox{note}{
    enhanced,
    breakable,
    title={Nota},
    fonttitle=\bfseries,
    colback=white,
    colframe=gray!50,
    boxrule=0.6pt,
    arc=2mm
}

% =========================================================================
% CONFIGURACIÓN FINAL (encabezados y pies)
% =========================================================================
\pagestyle{fancy}
\fancyhf{}

\fancyhead[L]{\nouppercase{\leftmark}}
\fancyhead[R]{\materia}

\fancyfoot[C]{\thepage}

\renewcommand{\headrulewidth}{0.4pt}
\renewcommand{\footrulewidth}{0pt}

\fancypagestyle{plain}{
    \fancyhf{}
    \fancyfoot[C]{\thepage}
    \renewcommand{\headrulewidth}{0pt}
}

% =========================================================================
% DOCUMENT
% =========================================================================
\begin{document}

% -------------------------------------------------------------------------
% PORTADA
% -------------------------------------------------------------------------
\begin{titlepage}

\thispagestyle{empty}

\begin{center}

\vspace*{1cm}

{\large \textsc{\universidad}}

\vspace{0.2cm}

{\large \carrera}

\vspace{3cm}

{\Huge \textbf{Apuntes de \materia}}

\vspace{1cm}

{\LARGE \titulo}

\vspace{0.8cm}

\textit{Comprender el proceso es el primer paso para transformar el estudio en una defensa sólida.}

\vspace{1.2cm}

\rule{0.70\textwidth}{0.5pt}

\vspace{0.5cm}

{\large Apuntes de estudio}

\vspace{0.5cm}

\rule{0.70\textwidth}{0.5pt}

\vfill

{\large \autor}

\vspace{0.3cm}

{\large \anio}

\end{center}

\vfill

\begin{flushleft}
\textit{Este es un modesto aporte a los estudiantes de ``Legislación y Práctica Profesional'' no pretende sustituir el material oficial de la materia.}
\end{flushleft}

\end{titlepage}

% -------------------------------------------------------------------------
% ÍNDICE
% -------------------------------------------------------------------------
\tableofcontents

% -------------------------------------------------------------------------
% INTRODUCCIÓN
% -------------------------------------------------------------------------
\chapter*{Presentación del tema}
\addcontentsline{toc}{chapter}{Presentación del tema}

Estos apuntes recorren el modo en que un hecho de la vida cotidiana se transforma en un conflicto jurídico y en la manera en que el abogado construye, a partir de él, una estrategia procesal completa: la demanda, el relato de los hechos y la prueba que lo sostiene, sabiendo siempre que la contraparte también despliega su propia estrategia.

A partir de ese núcleo se estudia en profundidad la prueba pericial (el perito, los puntos de pericia, el consultor técnico y la impugnación del dictamen), se analiza un caso práctico de accidente laboral y su tránsito por la vía administrativa y judicial, y se cierra con los demás medios de prueba relevantes para el litigio: la documental, la informativa y la testimonial, incluyendo los criterios para preparar a un testigo y formularle preguntas correctamente.

\begin{note}
El conocimiento del proceso judicial resulta relevante no solo para quien ejerce la abogacía, sino también para arquitectos e ingenieros, que suelen intervenir como peritos, consultores técnicos o testigos cuando una obra presenta fallas o filtraciones, así como para cualquier persona que deba enfrentar un conflicto y necesite saber qué documentos conservar, qué reconocer o desconocer, y a quién ofrecer como testigo.
\end{note}

Un mismo caso guía buena parte del desarrollo: una persona sufre un accidente automovilístico (``me chocaron el auto'') y, a partir de allí, se investiga el hecho, se evalúan las probabilidades de éxito, se diseña la demanda, se ofrece prueba pericial para tasar el daño, se compara con el procedimiento de un accidente laboral ante la Aseguradora de Riesgos del Trabajo (ART) y, finalmente, se prueban los hechos con documentos, informes y testigos.

% =========================================================================
% CAPÍTULO 1
% =========================================================================
\chapter{El conflicto y la estrategia procesal}

\section{El conflicto y sus vías de solución}

No todo hecho genera un conflicto, pero cuando esto ocurre, las partes cuentan con distintas vías para resolverlo. Algunas logran hacerlo por sí mismas, mediante un acuerdo amigable. Cuando esto no es posible, pueden recurrir a un mediador, a través de una mediación, o bien a un arbitraje, mecanismo todavía poco extendido en el ámbito local salvo en algunas empresas, aunque en expansión creciente.

\begin{important}
Cuando la justicia estatal no ofrece una respuesta suficiente en tiempo o en profundidad, no resulta desacertado recurrir a vías alternativas de resolución de conflictos como la mediación o el arbitraje. Cuando ninguna de estas vías resulta posible o suficiente, queda planteado el proceso judicial.
\end{important}

\section{El rol del abogado y la inteligencia natural}

\subsection{La especialidad del abogado es manejar el proceso}

La especialidad del abogado consiste, fundamentalmente, en el manejo del proceso judicial. El resto de los contenidos normativos de otras disciplinas suele ser conocido, muchas veces con mayor profundidad, por quienes ejercen o desarrollan las actividades reguladas por esas normas, no por ser abogados sino por tener que aplicarlas cotidianamente.

\begin{example}{Quienes conocen su materia mejor que el abogado}
Quienes se desempeñan en la actividad aeronáutica conocen en profundidad las normas de ese ámbito por tener que aplicarlas en su trabajo diario, y no por haber estudiado derecho. Lo mismo ocurre con cuestiones marítimas o con el derecho laboral: quienes participan de una negociación sindical con frecuencia dominan esa materia mejor que el propio abogado.
\end{example}

Lo que el abogado debe saber, entonces, es fundamentalmente el proceso: la herramienta que permite interpretar las normas cuando surge un problema y ponerlo en valor frente a un juez, con el fin de obtener un resultado favorable para el cliente.

\subsection{Qué significa ``resultado favorable''}

\begin{important}
Un resultado favorable no equivale necesariamente a ganar el cien por ciento del reclamo. Puede consistir en obtener una pérdida menor, una ganancia parcial o, en definitiva, ``obtener lo que se puede''. El abogado debe plantearse el objetivo del proceso en esos términos.
\end{important}

\subsection{Primero la inteligencia natural, después la artificial}

Antes de recurrir a la inteligencia artificial es necesario tener incorporada una base de inteligencia natural sobre el funcionamiento del proceso. De lo contrario, se corre el riesgo de quedar sujeto a lo que la herramienta devuelva, sin poder evaluar críticamente esa respuesta.

\begin{important}
Para poder formular una consulta adecuada a la inteligencia artificial es necesario saber previamente cómo funciona el proceso judicial. Si no se sabe qué preguntar, el riesgo de recibir una respuesta desacertada o inconducente es extraordinario.
\end{important}

\section{Del cliente a la estrategia}

\subsection{El conflicto llega a través del cliente}

La primera etapa del trabajo del abogado consiste en tomar conocimiento del conflicto a través del cliente, dado que el profesional no presenció los hechos. De allí la importancia de investigar adecuadamente y de formular la mayor cantidad de preguntas posibles para acercarse, en la medida en que sea posible, a lo efectivamente ocurrido.

\subsection{Evaluar el derecho y las probabilidades del caso}

Una vez reconstruido el hecho, corresponde evaluar el derecho aplicable y decidir si se recomienda o no iniciar el juicio. Es una obligación del abogado comunicarle al cliente una estimación de las probabilidades de éxito del caso, del mismo modo en que cualquier actividad requiere organizar previamente una estrategia.

\subsection{Estrategia personal, estrategia deportiva y estrategia procesal}

La diferencia central entre una estrategia personal (por ejemplo, para mejorar la condición física) y una estrategia procesal o competitiva radica en que, en esta última, el sujeto no actúa solo: debe tener en cuenta que otro también participa del juego, y en ocasiones lo hace muy bien. La estrategia procesal, por lo tanto, no puede pensarse únicamente en función de lo que la parte quiere, sino también de lo que se supone que la contraparte puede sostener y del modo en que ello puede afectar el resultado. Asimismo, exige un alto grado de flexibilidad, porque lo que ocurra no necesariamente coincidirá con lo planeado.

\begin{table}[H]
\centering
\caption{Estrategia personal frente a estrategia procesal}
\begin{tabularx}{\textwidth}{>{\raggedright\arraybackslash}p{0.32\textwidth} X X}
\toprule
 & \textbf{Estrategia personal (gimnasio, régimen)} & \textbf{Estrategia competitiva / procesal} \\
\midrule
\textbf{¿Con quién se enfrenta?} & Con uno mismo; no hay un rival que juegue. & Con otro que también juega, y a veces juega muy bien. \\
\textbf{¿Qué debe considerar?} & Lo que uno quiere y planifica. & Lo que uno quiere y lo que se supone que el otro puede sostener y afectará. \\
\textbf{¿Qué grado de flexibilidad exige?} & Bajo. & Alto: lo que ocurra no será necesariamente lo que uno quiere que pase. \\
\bottomrule
\end{tabularx}
\end{table}

\section{La demanda, el relato y las pruebas}

\subsection{La demanda como estrategia}

\begin{definition}{Demanda}
La demanda es la estrategia procesal del actor: el acto mediante el cual se fijan los hechos controvertidos a través de un relato, acompañado de las pruebas que lo sostienen o que se ofrecen para producirse durante el proceso.
\end{definition}

\subsection{El relato: por qué pensamos en historias}

Las personas se comunican fundamentalmente a través de relatos. El pensamiento humano necesita una estructura narrativa: una historia que empieza, se desarrolla y termina, en la medida en que se trata de seres temporales que no perciben la realidad de manera instantánea. Por ese motivo, la demanda —al igual que cualquier otro acto de comunicación jurídica relevante— se construye como un relato.

\subsection{Las pruebas y la estrategia de la contraparte}

Al relato de la demanda se le agregan las pruebas: aquellas que ya se poseen y se acompañan, y aquellas que se ofrecen para producirse durante el proceso, sin poder anticipar con exactitud su resultado. La parte actora supone que los testigos declararán en el sentido esperado y que los peritos se expedirán medianamente a favor de su posición, pero debe tener presente que la contraparte desplegará, a su vez, su propia estrategia.

\subsection{Verdad material y verdad legal formal}

Existe una diferencia entre el hecho realmente ocurrido —la verdad material, que ni siquiera puede conocerse con certeza absoluta— y la verdad legal formal, que es la conclusión a la que arriba el juez luego del período de prueba. Ambas verdades pueden coincidir en mayor o menor medida, pero no son necesariamente idénticas.

\begin{table}[H]
\centering
\caption{Verdad material, transcurso de la prueba y verdad legal formal}
\begin{tabularx}{\textwidth}{>{\raggedright\arraybackslash}X >{\raggedright\arraybackslash}X >{\raggedright\arraybackslash}X}
\toprule
\textbf{Verdad material} & \textbf{Transcurso de la prueba} & \textbf{Verdad legal formal} \\
\midrule
Lo que realmente ocurrió: la realidad, ``la verdad'', que ni siquiera se conoce con certeza. & Las pruebas se van produciendo; lo narrado se va probando. & Lo que el juez decide y termina estableciendo que ocurrió. Puede parecerse poco o mucho a la verdad material. \\
\bottomrule
\end{tabularx}
\end{table}

\begin{important}
El derecho acepta esta diferencia entre lo ocurrido y lo decidido porque, como ciencia social, no puede garantizar el acceso a la verdad plena. Aun así, puede perfeccionarse mediante mejores investigadores, mejores jueces y mejores abogados.
\end{important}

\subsection{La función social del abogado}

Más allá de esta limitación estructural, la función del abogado consiste en minimizar los conflictos y ofrecerles una solución, dado que el conflicto es inherente al ser humano y su resolución —aunque difícil y estresante— permite a las personas continuar con su vida.

\begin{example}{La empresa frente a un reclamo}
Una empresa a la que se le reclama una suma elevada evalúa si conviene resistir el reclamo o llegar a un acuerdo por una suma menor, con el fin de dar por concluido el conflicto y continuar con su actividad. La misma lógica se aplica a numerosas situaciones cotidianas: un despido, la compra de un inmueble, entre otras.
\end{example}

\section{El proceso como método}

\subsection{La estructura del proceso}

\begin{definition}{Estructura del proceso}
El proceso judicial se estructura sobre la base de tres posiciones: una parte actora que formula un reclamo o pedido; una parte demandada, a la que se le da traslado para que conteste; y un tercero —juez o árbitro— que resuelve el conflicto, no por ocupar una posición de superioridad, sino por el lugar que le corresponde en ese esquema.
\end{definition}

\subsection{El derecho no asegura justicia: garantiza un método}

\begin{important}
El derecho no puede garantizar la justicia del resultado. Lo que sí garantiza es un método o mecanismo detallado que asegure la igualdad de las partes ante el tercero que resuelve, junto con mecanismos para intentar convencer a ese tercero y, en caso de una resolución desfavorable e incorrecta, para recurrir a otra instancia.
\end{important}

\subsection{Estrategia de principio a fin}

Es fundamental sostener la estructura de pensamiento estratégico durante todo el proceso, desde su inicio —al demandar o al contestar la demanda— hasta su finalización. La ausencia de una estrategia definida equivale a afrontar el litigio sin rumbo.

\subsection{Por qué conocer el derecho procesal}

El conocimiento de las instituciones del derecho procesal resulta imprescindible incluso para quien domina otras ramas del derecho, porque frente a un conflicto concreto lo que se necesita, en primer lugar, es saber cómo son los mecanismos para reclamar y hacer efectivo un derecho.

\subsection{Cómo piensa cada rama del derecho}

Cada rama del derecho organiza el razonamiento del abogado de un modo distinto, según el tipo de cuestión que aborda.

\begin{table}[H]
\centering
\caption{Modo de razonar según la rama del derecho}
\begin{tabularx}{\textwidth}{>{\raggedright\arraybackslash}p{0.28\textwidth} X}
\toprule
\textbf{Rama del derecho} & \textbf{Cómo se piensa} \\
\midrule
\textbf{Derecho civil (daños)} & Las cosas eran de un modo, se produjo un cambio y ahora son de otro; la diferencia entre lo que era y lo que es constituye el daño. Luego corresponde analizar cómo se cuantifica y cómo se prueba. \\
\textbf{Derecho administrativo} & Desde la posición del administrado: la administración puede resultar arbitraria al decidir; existen principios que debió respetar necesariamente. Si no los respetó y ello causó un daño, eso es lo que se reclama. \\
\textbf{Derecho penal} & Corresponde analizar si el hecho constituye o no un delito y si la persona a quien se le atribuye es o no responsable. \\
\bottomrule
\end{tabularx}
\end{table}

\section*{Cornell: el conflicto y la estrategia procesal}
\addcontentsline{toc}{section}{Cornell: el conflicto y la estrategia procesal}

\begin{table}[H]
\centering
\begin{tabularx}{\textwidth}{
    >{\raggedright\arraybackslash}p{0.30\textwidth}
    >{\raggedright\arraybackslash}X
}
\toprule
\textbf{Preguntas / palabras clave} & \textbf{Notas / respuestas} \\
\midrule

\textbf{¿Qué vías existen para resolver un conflicto?} &
El acuerdo directo entre las partes, la mediación, el arbitraje y, cuando estas vías no son posibles o suficientes, el proceso judicial. \\

\textbf{¿Cuál es la especialidad propia del abogado?} &
Manejar el proceso, de modo de poder interpretar las normas de cualquier rama del derecho y ponerlas en valor ante un juez para obtener un resultado favorable, que no siempre equivale a ganar por completo. \\

\textbf{¿Por qué es necesario saber cómo funciona el proceso antes de usar inteligencia artificial?} &
Porque la inteligencia artificial devuelve respuestas de baja calidad cuando no se sabe qué preguntarle; contar con una base de inteligencia natural evita depender ciegamente de esas respuestas. \\

\textbf{¿Qué diferencia a la estrategia procesal de una estrategia personal?} &
En el proceso la contraparte también despliega su propia estrategia, por lo que es necesario prever lo que puede sostener y actuar con un alto grado de flexibilidad. \\

\textbf{¿Qué es la demanda y qué papel cumple el relato?} &
La demanda es la estrategia del actor: fija los hechos controvertidos mediante un relato y los acompaña con la prueba correspondiente. \\

\textbf{¿En qué se diferencian la verdad material y la verdad legal formal?} &
La verdad material es lo que realmente ocurrió, que no puede conocerse con certeza absoluta; la verdad legal formal es la conclusión a la que arriba el juez tras la prueba, y puede diferir de la primera. \\

\textbf{¿Qué garantiza el derecho si no puede asegurar la justicia?} &
Garantiza un método que asegura la igualdad de las partes ante el tercero que resuelve, con mecanismos para convencerlo y para recurrir ante una decisión incorrecta. \\

\midrule
\multicolumn{2}{p{\textwidth}}{
\textbf{Resumen:} El conflicto, cuando no puede resolverse entre las partes ni por mediación o arbitraje, se somete a un proceso judicial cuya función no es garantizar la justicia sino un método igualitario. El abogado debe manejar ese método y traducir el conflicto en una demanda que sea, en sí misma, una estrategia: un relato de los hechos sostenido por pruebas, elaborado sabiendo que la contraparte también actúa estratégicamente y que la verdad que finalmente se establece en el proceso puede no coincidir exactamente con lo realmente ocurrido.
} \\
\bottomrule
\end{tabularx}
\end{table}

% =========================================================================
% CAPÍTULO 2
% =========================================================================
\chapter{La prueba pericial}

\section{Régimen de la prueba pericial}

\subsection{Procedencia}

La prueba pericial resulta necesaria cuando, para apreciar un hecho controvertido dentro del proceso, se requiere un conocimiento técnico que excede al del juez. Para suplir esa carencia interviene un auxiliar de la justicia: el perito.

\begin{formula}{Código Procesal Civil y Comercial de la Nación, art. 457 --- Procedencia}
Será admisible la prueba pericial cuando la apreciación de los hechos controvertidos requiera conocimientos especiales en alguna ciencia, arte, industria o actividad técnica especializada.
\end{formula}

\subsection{Perito único y consultores técnicos}

La prueba pericial está a cargo, por regla general, de un único perito, designado de oficio por el juez. Existen dos excepciones a esta regla: los juicios de inhabilitación e incapacidad y los juicios de nulidad testamentaria.

\begin{formula}{CPCCN, art. 458 --- Perito único. Consultores técnicos}
La prueba pericial estará a cargo de UN (1) perito único designado de oficio por el juez, salvo cuando una ley especial establezca un régimen distinto. En los procesos de declaración de incapacidad y de inhabilitación, se estará a lo dispuesto en el artículo 626, inciso 3. En el juicio por nulidad de testamento, el juez podrá nombrar de oficio TRES (3) peritos cuando por la importancia y complejidad del asunto lo considere conveniente. Si los peritos fuesen TRES (3), el juez les impartirá las directivas sobre el modo de proceder para realizar las operaciones tendientes a la producción y presentación del dictamen. Cada parte tiene la facultad de designar un consultor técnico.
\end{formula}

\subsection{Ofrecimiento: especialización y puntos de pericia}

\begin{definition}{Puntos de pericia}
Los puntos de pericia son los temas concretos sobre los cuales el perito debe trabajar y expedir su informe. Deben indicarse al momento de ofrecer la prueba, junto con la especialización que debe tener el perito.
\end{definition}

\begin{formula}{CPCCN, art. 459 --- Designación. Puntos de pericia (primer párrafo)}
Al ofrecer la prueba pericial se indicará la especialización que ha de tener el perito y se propondrán los puntos de pericia; si la parte ejerciera la facultad de designar consultor técnico, deberá indicar, en el mismo escrito, su nombre, profesión y domicilio.
\end{formula}

Una vez corrido el traslado del ofrecimiento de esta prueba, la contraparte puede impugnarlo, añadir puntos de pericia u observar los ya propuestos.

\begin{note}
Si las partes no acuerdan el perito y los puntos de pericia, el juez los fija de oficio en la denominada ``audiencia del artículo 360'' del CPCCN, pudiendo agregar o eliminar puntos. El plazo para presentar el dictamen, si nada se dice, es por principio general de quince días.
\end{note}

\section{Título, recusación, aceptación y remoción del perito}

\begin{table}[H]
\centering
\caption{Reglas sobre la actuación del perito}
\begin{tabularx}{\textwidth}{>{\raggedright\arraybackslash}p{0.25\textwidth} X}
\toprule
\textbf{Tema} & \textbf{Regla} \\
\midrule
\textbf{Título habilitante} & Si la actividad está reglamentada, el perito debe presentar título habilitante. Si no existe perito con título en la materia, puede designarse a otra persona con conocimientos suficientes. \\
\textbf{Recusación} & Pueden recusarlo las partes, invocando causa justa, o el juez de oficio, por falta de título o incompetencia. Si el perito reconoce la causa o guarda silencio, se lo reemplaza; si la niega, la cuestión tramita por separado, sin interrumpir el proceso. \\
\textbf{Aceptación del cargo} & Debe aceptarlo dentro del tercer día; si carece de título habilitante, debe jurar desempeñar fielmente el cargo. \\
\textbf{Remoción} & El juez puede removerlo de oficio si renuncia, si no presenta el informe dentro del plazo o si se rehúsa a entregarlo. \\
\bottomrule
\end{tabularx}
\end{table}

En la práctica actual, la aceptación del cargo suele formalizarse mediante un escrito presentado por vía electrónica, en el que el perito manifiesta que acepta el cargo y jura desempeñarlo, sin necesidad de comparecer personalmente a firmar como se exigía anteriormente.

\section{Adelanto de gastos, idóneos y consultor técnico}

\subsection{El adelanto de gastos y sus consecuencias}

El perito puede solicitar que se le adelanten los gastos necesarios para la realización de la pericia, y el juez resuelve sobre ese pedido. Las partes pueden oponerse u observar el monto solicitado antes de que el juez decida.

\begin{important}
Si el juez concede el adelanto de gastos pedido por el perito y la parte obligada no lo cumple, la prueba pericial no se produce. La parte se queda, entonces, sin una prueba que puede ser fundamental para su caso.
\end{important}

\begin{note}
El beneficio de litigar sin gastos permite, al menos temporariamente, no afrontar los gastos del proceso, entre ellos los de la pericia.
% TODO: contenido incompleto en las notas originales; el docente anunció que este beneficio se trataría en detalle más adelante en el curso, pero el desarrollo no figura en este apunte.
\end{note}

\subsection{Los peritos y los idóneos}

Los peritos se inscriben anualmente en las listas de la cámara correspondiente (civil, penal, etcétera). Cuando, excepcionalmente, no existiera un perito matriculado de determinada especialidad, puede llamarse a una persona idónea en la materia.

\begin{important}
La prueba pericial es un tema especialmente sensible: existen numerosas cuestiones técnicas —por ejemplo, si un edificio estuvo bien construido o si hubo una falla estructural, o si existió una mala praxis médica— que exceden el conocimiento del abogado, por lo que buena parte de la labor profesional consiste en trabajar adecuadamente con la pericia, se cuente o no con un consultor técnico propio.
\end{important}

\subsection{El consultor técnico}

\begin{table}[H]
\centering
\caption{Perito frente a consultor técnico}
\begin{tabularx}{\textwidth}{>{\raggedright\arraybackslash}p{0.22\textwidth} X X}
\toprule
 & \textbf{Perito} & \textbf{Consultor técnico} \\
\midrule
\textbf{¿Quién lo designa?} & El juez, de oficio (perito único). & Cada parte tiene la facultad de designar uno. \\
\textbf{¿Presenta dictamen?} & Sí, presenta el informe pericial. & No presenta dictamen; no puede actuar como perito oficial. \\
\textbf{¿Qué hace?} & Acepta el cargo, realiza las operaciones periciales y presenta el informe dentro del plazo. & Asiste a las operaciones periciales, acompaña al perito, puede sugerirle aspectos a considerar y asesora a la parte para consentir u objetar la pericia. \\
\textbf{Requisitos} & Título habilitante si la actividad está reglamentada; inscripción anual en la cámara; puede ser recusado o removido. & Es un profesional o idóneo de confianza de la parte. \\
\bottomrule
\end{tabularx}
\end{table}

\begin{example}{Las filtraciones en un departamento}
Ante una demanda por deterioros o filtraciones en un edificio, el perito oficial debe determinar el origen del problema, cómo resolverlo y su costo. El consultor técnico de parte puede asistir a las operaciones periciales, acompañar al perito en el momento en que se toman fotografías y sugerirle, por ejemplo, que registre otro sector o que analice determinado aspecto adicional.
\end{example}

\subsection{Impugnar y pedir aclaraciones}

\begin{important}
La objeción técnica de la pericia se denomina impugnación. Resulta tan importante poder impugnar el dictamen como poder solicitar aclaraciones sobre él. El diseño de los puntos de pericia es decisivo, dado que en la práctica el juez suele otorgarle a la pericia un valor casi determinante.
\end{important}

\section{Remoción, impugnación y valoración de la pericia}

\subsection{Las formas de pedir la remoción del perito}

Existen tres motivos por los cuales puede solicitarse la remoción del perito: que no presente el informe o dictamen; que se niegue a hacerlo; o que renuncie sin un motivo atendible. Un motivo atendible sería, por ejemplo, una enfermedad sobreviniente que le impida continuar, supuesto en el cual puede concederse una prórroga.

\subsection{Ejemplo: la pericia médica en accidentes laborales}

\begin{example}{Recorrido de la pericia médica en un accidente laboral}
\begin{enumerate}
\item Se sortea el perito médico, quien debe aceptar el cargo.
\item El perito cita a la persona, la revisa y le solicita los estudios complementarios que considere necesarios (por ejemplo, radiografías).
\item Con la información recibida, elabora el informe dentro del plazo fijado; si no lo presenta, puede ser removido.
\item Se da conocimiento del informe a ambas partes, que pueden impugnarlo o solicitar aclaraciones.
\item El juez lo tiene en cuenta al dictar sentencia.
\end{enumerate}
\end{example}

\subsection{Impugnar la pericia cuando el perito ``no encuentra nada''}

Cuando el perito concluye que el trabajador no presenta incapacidad y la parte actora sostiene lo contrario, corresponde impugnar el dictamen señalando con precisión sus deficiencias: la falta de aplicación de un método determinado, la omisión de un estudio necesario o la falta de consideración de algún elemento relevante. Según el caso, puede ordenarse una nueva pericia o rechazarse la impugnación.

\subsection{Valoración de la pericia por el juez}

\begin{important}
El juez tiene en cuenta el informe pericial al dictar sentencia, pero no está obligado a basarse necesariamente en él: puede acogerlo o apartarse de sus conclusiones, siempre que cuente con fundamento para hacerlo. De allí la importancia de impugnar y solicitar aclaraciones en tiempo oportuno.
\end{important}

\section*{Cornell: la prueba pericial}
\addcontentsline{toc}{section}{Cornell: la prueba pericial}

\begin{table}[H]
\centering
\begin{tabularx}{\textwidth}{
    >{\raggedright\arraybackslash}p{0.30\textwidth}
    >{\raggedright\arraybackslash}X
}
\toprule
\textbf{Preguntas / palabras clave} & \textbf{Notas / respuestas} \\
\midrule

\textbf{¿Cuándo procede la prueba pericial?} &
Cuando la apreciación de un hecho controvertido requiere conocimientos especiales en alguna ciencia, arte, industria o actividad técnica especializada (art. 457 CPCCN). \\

\textbf{¿Quién designa al perito y qué excepciones existen a la regla del perito único?} &
Lo designa de oficio el juez, como perito único, salvo en los juicios de inhabilitación e incapacidad y en los juicios de nulidad testamentaria, donde pueden nombrarse tres peritos. \\

\textbf{¿Qué son los puntos de pericia y quién puede observarlos?} &
Son los temas sobre los que el perito debe expedirse; se proponen al ofrecer la prueba y la contraparte puede impugnarlos, añadir puntos u observarlos. Si no hay acuerdo, el juez los fija en la audiencia del artículo 360. \\

\textbf{¿Qué ocurre si la parte no cumple el adelanto de gastos concedido al perito?} &
La prueba pericial no se produce, y la parte se queda sin una prueba que puede resultar fundamental para su caso. \\

\textbf{¿En qué se diferencia el consultor técnico del perito?} &
El perito es designado de oficio y presenta el dictamen oficial; el consultor técnico es designado por la parte, no presenta dictamen, sino que asiste a las operaciones periciales y asesora a la parte. \\

\textbf{¿Cuáles son las causales de remoción del perito?} &
No presentar el informe en plazo, negarse a presentarlo, o renunciar sin motivo atendible. \\

\textbf{¿Qué puede hacer una parte frente a un dictamen desfavorable?} &
Impugnarlo, señalando con precisión sus deficiencias, o solicitar aclaraciones; el juez puede ordenar una nueva pericia. \\

\textbf{¿Está el juez obligado a seguir el dictamen pericial?} &
No; puede apartarse de sus conclusiones, pero necesita fundamento para hacerlo. \\

\midrule
\multicolumn{2}{p{\textwidth}}{
\textbf{Resumen:} La prueba pericial suple el conocimiento técnico que excede al juez y, por regla, la produce un único perito de oficio. Su eficacia depende de un correcto diseño de los puntos de pericia, de la atención al adelanto de gastos —cuyo incumplimiento puede dejar a la parte sin prueba— y del rol activo del consultor técnico, que asesora sin sustituir al perito. El resultado del dictamen no vincula automáticamente al juez, que puede apartarse de él con fundamento, razón por la cual impugnar y pedir aclaraciones resulta decisivo.
} \\
\bottomrule
\end{tabularx}
\end{table}

% =========================================================================
% CAPÍTULO 3
% =========================================================================
\chapter{Caso práctico: el accidente laboral}

\section{La ART y la revisión judicial}

\begin{definition}{Accidente laboral}
Es aquel que se produce con ocasión del trabajo o durante el trayecto que la persona realiza entre su domicilio y el lugar de trabajo, en cualquiera de los dos sentidos.
\end{definition}

\begin{example}{El trabajador del balancín}
Un trabajador se accidenta con un balancín y se corta un dedo. Acude a la Comisión Médica, órgano de la Superintendencia de Riesgos del Trabajo, para que un médico lo revise y fije un porcentaje de incapacidad. Ese porcentaje suele ser reducido o directamente no reconocido, motivo por el cual el trabajador inicia un juicio contra la ART reclamando la indemnización que considera que le corresponde.
% TODO: en el apunte original se menciona a la "Superintendencia del Trabajo" como institución que revisa al trabajador; se trata, en rigor, de la Comisión Médica dependiente de la Superintendencia de Riesgos del Trabajo, según aclaración incorporada por quien transcribió la clase.
\end{example}

\begin{important}
La decisión de la ART, por tratarse de un acto de carácter administrativo, no cierra el conflicto. Toda decisión administrativa, para ser válida, debe poder revisarse en sede judicial; de lo contrario, se privaría a la persona de una verdadera defensa en juicio. El plazo para impugnar esa decisión mediante una demanda suele ser breve: si no se ejerce dentro de ese plazo, la resolución administrativa queda firme, aunque la posibilidad de impugnarla existió.
\end{important}

Frente a una resolución desfavorable de la Comisión Médica, corresponde esperar a que quede formalizada y, dentro del plazo previsto, optar por la vía de apelación o por el inicio directo del juicio.
% TODO: en el apunte original el alumno menciona un plazo de apelación de quince días, y luego duda y menciona alternativamente dos días; el dato no quedó precisado con certeza en la clase. Se indica, asimismo, que tanto la vía de apelación como el juicio ordinario insumen un plazo similar de inicio, de aproximadamente siete u ocho meses.

\section{Audiencias y conciliación}

En el fuero laboral existen audiencias de conciliación, distintas de la audiencia preliminar del proceso civil. Algunos jueces convocan a las partes a esa audiencia aun cuando no se los haya solicitado expresamente.

\begin{important}
En la audiencia de conciliación, el juez que ya conoce el caso y las pruebas producidas puede orientar a ambas partes hacia un acuerdo, sin adelantar el sentido de su futura decisión, pero mostrándoles hacia dónde se dirige previsiblemente el proceso. De ese modo, con cierta presión moderada, puede lograr que las partes concilien, en el entendimiento de que un acuerdo suele resultar más económico, más sencillo y más rápido que continuar el litigio hasta la sentencia.
\end{important}

\section*{Cornell: el caso del accidente laboral}
\addcontentsline{toc}{section}{Cornell: el caso del accidente laboral}

\begin{table}[H]
\centering
\begin{tabularx}{\textwidth}{
    >{\raggedright\arraybackslash}p{0.30\textwidth}
    >{\raggedright\arraybackslash}X
}
\toprule
\textbf{Preguntas / palabras clave} & \textbf{Notas / respuestas} \\
\midrule

\textbf{¿Qué se considera un accidente laboral?} &
El que se produce con ocasión del trabajo o durante el trayecto de ida o de vuelta entre el domicilio y el lugar de trabajo. \\

\textbf{¿Por qué la decisión de la ART no cierra el conflicto?} &
Porque se trata de un órgano administrativo, y toda decisión administrativa debe poder revisarse judicialmente para garantizar la defensa en juicio; de no impugnarse dentro del plazo, la resolución queda firme. \\

\textbf{¿Qué función cumple la audiencia de conciliación en el fuero laboral?} &
Permite al juez, sin adelantar su decisión, orientar a las partes hacia un acuerdo que suele resultar más económico, más simple y más rápido que continuar el proceso. \\

\midrule
\multicolumn{2}{p{\textwidth}}{
\textbf{Resumen:} El caso del accidente laboral muestra en la práctica que una decisión administrativa —como la de la ART sobre el porcentaje de incapacidad— nunca cierra definitivamente el conflicto, porque siempre existe la posibilidad de someterla a revisión judicial dentro de un plazo determinado. Iniciado el proceso judicial, la audiencia de conciliación propia del fuero laboral ofrece a las partes una vía para llegar a un acuerdo antes de la sentencia.
} \\
\bottomrule
\end{tabularx}
\end{table}

% =========================================================================
% CAPÍTULO 4
% =========================================================================
\chapter{Documental, informativa y testimonial}

\section{El testigo y qué se quiere acreditar}

\begin{definition}{Testigo}
Es una persona que, a través de sus sentidos, percibió hechos ocurridos en el momento relevante para el proceso y que se ofrece para que los acredite o declare sobre ellos.
\end{definition}

No todo testigo declara sobre el hecho principal: también puede ofrecerse un testigo para acreditar sus consecuencias, como ocurre con quienes conocieron la vida de una persona antes y después de un accidente.

\begin{table}[H]
\centering
\caption{Testigos del hecho y testigos de las consecuencias}
\begin{tabularx}{\textwidth}{>{\raggedright\arraybackslash}p{0.22\textwidth} X X}
\toprule
 & \textbf{Testigo del hecho} & \textbf{Testigo de las consecuencias} \\
\midrule
\textbf{¿Quién podría serlo?} & Una persona presente en el momento del hecho (por ejemplo, quien presenció un accidente). & Amigos o conocidos del actor. \\
\textbf{¿Qué acredita?} & El hecho mismo: lo ocurrido en el momento. & La situación anterior y posterior de la persona, que suele emplearse para acreditar el daño moral. \\
\bottomrule
\end{tabularx}
\end{table}

\begin{example}{El accidente en el barco}
En el caso de referencia, el actor sufrió un accidente náutico en el que perdió a su esposa. Se busca acreditar: el hecho mismo, mediante testigos presenciales; el daño moral, con amigos y conocidos que dan cuenta de su vida antes y después del accidente; y los ingresos que dejó de percibir, mediante prueba informativa.
\end{example}

\section{Medios de prueba: la confesión y la documental}

Entre los medios para acreditar los hechos se encuentra la confesión, de utilidad limitada, y la prueba documental, que resulta fundamental.

\begin{important}
La prueba documental debe ofrecerse junto con la demanda. Si el documento no se posee al momento de presentarla, igualmente se cumple con la carga procesal indicando cuál es el documento, dónde se encuentra y todos los elementos disponibles sobre él —como ocurre, por ejemplo, con una historia clínica que obra en un hospital.
\end{important}

\begin{note}
No es posible presentar documentos no ofrecidos con la demanda en la audiencia del artículo 360, salvo que se trate de un hecho nuevo, supuesto en el cual corresponde iniciar el incidente respectivo mediante un escrito específico.
\end{note}

La demanda debe contener el relato de los hechos, el derecho invocado, los daños reclamados, los fundamentos de la responsabilidad, el objeto de la pretensión y el ofrecimiento de prueba, acompañada de toda la documental disponible. Al ser notificado, el demandado recibe la cédula junto con la copia de la demanda y de la documental, salvo que esta última sea muy voluminosa, en cuyo caso puede solicitarse que no se traslade físicamente por encontrarse disponible en el expediente digital.

\section{La carga del demandado ante los documentos}

Al contestar la demanda, el demandado tiene la carga de pronunciarse sobre cada documento acompañado: reconocerlo o desconocerlo. La omisión de pronunciarse implica su reconocimiento.

Frente a un documento que se pretende desconocer, corresponde hacerlo de manera puntual, indicando el motivo (por ejemplo, por inexacto o por falso), sin recurrir a un desconocimiento genérico de toda la prueba documental.

\begin{table}[H]
\centering
\caption{Instrumento público frente a instrumento privado}
\begin{tabularx}{\textwidth}{>{\raggedright\arraybackslash}p{0.22\textwidth} X X}
\toprule
 & \textbf{Instrumento público} & \textbf{Instrumento privado} \\
\midrule
\textbf{¿Cómo se lo cuestiona?} & Mediante redargución de falsedad, un procedimiento específico. & Se lo desconoce. \\
\textbf{¿Qué ocurre después?} & Se presume que hace fe pública; cuestionarlo es un procedimiento dificultoso. & Quien lo presentó debe probar que la firma es auténtica, o acreditarlo por otro medio. \\
\bottomrule
\end{tabularx}
\end{table}

\begin{important}
Desconocer sistemáticamente todos los documentos, incluso los que resultan evidentes, no constituye una estrategia procesal inteligente: le resta credibilidad a la parte frente al juez. Reconocer lo que no puede negarse y controvertir con precisión lo restante otorga mayor seriedad y contundencia a la posición procesal.
\end{important}

\section{Prueba informativa}

La prueba informativa consiste en solicitar, mediante un oficio, que una entidad pública o privada informe sobre determinada cuestión (por ejemplo, un oficio a un banco para que informe sobre una cuenta).

\begin{example}{Acreditar ingresos mediante oficios}
Para acreditar que una persona se desempeñaba como docente en determinadas universidades, puede libraserse un oficio a esas instituciones para que informen desde cuándo ejercía el cargo y cuánto percibía por él. De ese modo se acredita, indirectamente, el ingreso que la persona dejó de percibir a raíz del hecho dañoso, como parte de la estrategia probatoria.
\end{example}

\section{Preparación del testigo y generales de la ley}

\subsection{Por qué hay que preparar al testigo}

El testigo suele comparecer con temor. Prepararlo —sin que ello implique aleccionarlo— consiste en explicarle el procedimiento con claridad, de modo que pueda declarar con tranquilidad.

\begin{table}[H]
\centering
\caption{Modalidades de la prueba testimonial}
\begin{tabularx}{\textwidth}{>{\raggedright\arraybackslash}p{0.30\textwidth} X}
\toprule
\textbf{Modalidad} & \textbf{Cómo se desarrolla} \\
\midrule
\textbf{El tribunal interroga primero} & El juez o el secretario formulan todas las preguntas y luego consultan al abogado si desea ampliar el interrogatorio. \\
\textbf{Las partes formulan las preguntas} & El tribunal no interroga primero; son los abogados quienes formulan directamente las preguntas al testigo. \\
\bottomrule
\end{tabularx}
\end{table}

\begin{important}
Cuando un testigo declara deficientemente, en la mayoría de los casos la responsabilidad corresponde al abogado que lo ofreció, por no haber conversado previamente con él ni haberlo preparado adecuadamente.
\end{important}

\subsection{Las generales de la ley}

\begin{definition}{Generales de la ley}
Es el procedimiento inicial mediante el cual se interroga al testigo sobre su situación personal frente a las partes, previo a que declare sobre los hechos.
\end{definition}

\begin{important}
Al preparar al testigo corresponde explicarle: que declara bajo juramento y que, si miente, existen consecuencias penales, aunque ello forma parte normal del procedimiento y no debe generarle temor; que se le preguntarán datos personales; que se le consultará si conoce a la parte o mantiene algún vínculo con ella, debiendo responder con precisión y sinceridad sobre el grado de esa relación, sin que la amistad le impida ser objetivo; y que también se le preguntará si existen deudas entre él y la parte, o si es acreedor.
\end{important}

\section{Taller: cómo formular preguntas al testigo}

\subsection{La regla: la pregunta no debe conducir a la respuesta}

\begin{important}
La pregunta dirigida a un testigo nunca debe poder contestarse simplemente por sí o por no, ni debe presuponer en su formulación el dato que se pretende obtener. Corresponde comenzar con preguntas amplias y abiertas, y ir avanzando sobre la base de las respuestas que el testigo brinda. Frente a una pregunta inductiva formulada por la contraparte o por el tribunal, corresponde oponerse.
\end{important}

\begin{table}[H]
\centering
\caption{Preguntas cerradas o inductivas frente a preguntas abiertas}
\begin{tabularx}{\textwidth}{>{\raggedright\arraybackslash}X >{\raggedright\arraybackslash}X}
\toprule
\textbf{Pregunta cerrada o inductiva (a evitar)} & \textbf{Pregunta abierta (correcta)} \\
\midrule
``¿Hace cuánto formaban parte del club?'': presupone que el testigo conocía al actor y que ambos eran socios de un club. & ``¿De dónde lo conoce?'' \\
``¿Frecuentaba tal club?'': incorpora el dato del club en la propia pregunta, de modo que se responde por sí o por no. & ``¿Qué actividades realizaba el actor?'' \\
Cualquier pregunta que se conteste por sí o por no. & ``¿Con qué frecuencia se veían?'', ``¿dónde las realizaba?'', ``¿desde cuándo?'': cada pregunta se apoya en la respuesta anterior. \\
\bottomrule
\end{tabularx}
\end{table}

\begin{example}{El juez que preguntaba de manera indicativa}
En un juicio, el juez interrogó a un testigo con preguntas de tipo indicativo (por ejemplo, sugiriendo directamente la respuesta esperada). Al advertírsele que estaba formulando la pregunta de manera indicativa y solicitársele que preguntara, en cambio, qué actividad realizaba el testigo, el juez respondió que siempre había preguntado de ese modo, lo que derivó en un cuestionamiento posterior de lo actuado ante la Cámara.
% TODO: el audio original presenta fragmentos poco claros en esta anécdota; se conserva el sentido general expresado por el docente.
\end{example}

\subsection{Derechos de las partes y del tribunal al interrogar}

\begin{table}[H]
\centering
\caption{Orden y derechos en el interrogatorio del testigo}
\begin{tabularx}{\textwidth}{>{\raggedright\arraybackslash}p{0.22\textwidth} >{\raggedright\arraybackslash}p{0.28\textwidth} X}
\toprule
\textbf{Momento} & \textbf{Quién pregunta} & \textbf{Qué ocurre} \\
\midrule
Primero & La parte que ofreció al testigo (o el tribunal, según la modalidad). & Formula sus preguntas; el tribunal puede preguntar lo que considere que falta. \\
Luego & El tribunal consulta a la parte oferente. & La parte puede ampliar el interrogatorio o considerar que lo preguntado es suficiente. \\
Después & La contraparte, que no ofreció al testigo. & Tiene derecho a formular preguntas al testigo, aunque no haya sido su testigo. \\
Si hay discusión sobre las preguntas & --- & Se hace salir al testigo de la sala, se discute la cuestión entre las partes y el tribunal, y luego se lo hace volver a pasar. \\
\bottomrule
\end{tabularx}
\end{table}

\begin{note}
Quedaron anunciados como temas pendientes, no desarrollados en esta clase: la prueba confesional y el beneficio de litigar sin gastos.
\end{note}

\section*{Cornell: documental, informativa y testimonial}
\addcontentsline{toc}{section}{Cornell: documental, informativa y testimonial}

\begin{table}[H]
\centering
\begin{tabularx}{\textwidth}{
    >{\raggedright\arraybackslash}p{0.30\textwidth}
    >{\raggedright\arraybackslash}X
}
\toprule
\textbf{Preguntas / palabras clave} & \textbf{Notas / respuestas} \\
\midrule

\textbf{¿Qué diferencia hay entre un testigo del hecho y un testigo de las consecuencias?} &
El testigo del hecho acredita lo ocurrido en el momento; el testigo de las consecuencias (por ejemplo, amigos o conocidos) acredita la situación anterior y posterior, útil para probar el daño moral. \\

\textbf{¿Cuándo debe ofrecerse la prueba documental?} &
Con la demanda; si no se posee el documento en ese momento, se cumple indicando cuál es y dónde se encuentra. Solo puede incorporarse documental posterior mediante el incidente de hecho nuevo. \\

\textbf{¿Qué carga tiene el demandado frente a los documentos acompañados?} &
Debe reconocerlos o desconocerlos de manera puntual; el silencio implica reconocimiento. \\

\textbf{¿Cómo se cuestiona un instrumento público y cómo uno privado?} &
El instrumento público se cuestiona mediante redargución de falsedad; el instrumento privado simplemente se desconoce, y quien lo presentó debe probar su autenticidad. \\

\textbf{¿Cómo se acredita un dato mediante prueba informativa?} &
Mediante un oficio a una entidad pública o privada que informe sobre la cuestión de que se trate (por ejemplo, ingresos o cargos desempeñados). \\

\textbf{¿Qué hay que explicarle a un testigo antes de que declare?} &
Que declara bajo juramento, que existen consecuencias penales si miente aunque ello no debe generarle temor, que se le preguntarán datos personales y su relación con la parte, y si existen deudas. \\

\textbf{¿Cómo debe formularse una pregunta a un testigo?} &
De manera abierta, sin presuponer la respuesta ni poder contestarse por sí o por no; se avanza progresivamente sobre la base de lo que el testigo va respondiendo. \\

\midrule
\multicolumn{2}{p{\textwidth}}{
\textbf{Resumen:} Los hechos del proceso se acreditan mediante documentos ofrecidos con la demanda —que el demandado debe reconocer o desconocer con precisión—, mediante informes obtenidos a través de oficios y mediante testigos, que pueden declarar sobre el hecho mismo o sobre sus consecuencias. El testigo debe ser preparado explicándole el procedimiento y las generales de la ley, y debe ser interrogado con preguntas abiertas que no conduzcan la respuesta, respetando el orden y los derechos de cada parte durante la audiencia.
} \\
\bottomrule
\end{tabularx}
\end{table}

% =========================================================================
% SÍNTESIS GENERAL
% =========================================================================
\chapter*{Síntesis general}
\addcontentsline{toc}{chapter}{Síntesis general}

Todo conflicto que las partes no logran resolver por sí mismas termina, eventualmente, en un proceso judicial, cuya función social consiste en ofrecer una solución que permita a las personas continuar con su vida. El proceso no garantiza la justicia del resultado, sino un método que asegura la igualdad de las partes ante el tercero que resuelve; por ese motivo, la demanda debe concebirse como una estrategia flexible que fija un relato y ofrece prueba, sabiendo que la contraparte también actúa estratégicamente y que la verdad legal formal puede diferir de la verdad material.

La prueba pericial exige un perito único, puntos de pericia bien diseñados, atención al adelanto de gastos y, cuando sea posible, un consultor técnico que asesore a la parte; su resultado siempre puede impugnarse, porque el juez la valora sin quedar atado a sus conclusiones. El caso del accidente laboral muestra que lo decidido por un órgano administrativo es siempre revisable judicialmente, y que las audiencias de conciliación pueden ahorrar tiempo y costos a ambas partes. Finalmente, los hechos se acreditan con documentos ofrecidos junto con la demanda, con informes obtenidos mediante oficios y con testigos —del hecho y de sus consecuencias—, a quienes es necesario preparar adecuadamente y a quienes se interroga mediante preguntas abiertas que no conduzcan la respuesta.

\end{document}
